\chapter{Fundamentação Teórica}
\label{cap:fundamentacao-teorica}

Para entender como a união de diferentes disciplinas ajuda a criar soluções de
software melhores e mais eficientes, é preciso apresentar a base teórica desta
pesquisa. Este capítulo discute a mudança de um modelo de ensino dividido em
matérias isoladas para uma visão integrada, a importância das metodologias
ativas para dar protagonismo ao estudante e como isso acontece na prática por
meio do Projeto Integrador no curso de Engenharia de Software.

A assertividade no desenvolvimento de software — ou seja, a capacidade de
entregar exatamente o que é necessário da melhor forma possível — não depende
apenas de saber programar. Ela depende da capacidade de fazer as diferentes
áreas "conversarem" entre si. \citeonline{silva2024processo} destacam que,
quando as barreiras entre os conhecimentos são derrubadas, as soluções criadas
se tornam muito mais funcionais e adequadas.

Nesse cenário, o protagonismo do aluno é o que garante que a integração
funcione. Como afirma \citeonline{ribeiro2025interdisciplinaridade}, quando o
estudante assume o comando do seu aprendizado e usa a interdisciplinaridade
como ferramenta, ele desenvolve um pensamento crítico que permite criar
soluções que realmente funcionam no mundo real. Portanto, integrar matérias
como banco de dados, arquitetura e testes de forma simultânea não é apenas um
método de ensino, mas a estratégia central para formar profissionais capazes de
entregar tecnologia de alta qualidade e com valor real para a sociedade.

\section{Do Ensino Dividido à Integração de Saberes: Multidisciplinaridade e Interdisciplinaridade}\label{sec:do-ensino-dividido}

Historicamente, o ensino superior foi organizado de forma muito fragmentada,
separando o conhecimento em áreas isoladas. Segundo
\citeonline{ribeiro2025interdisciplinaridade}, o ensino ainda é fortemente
dividido em disciplinas isoladas, o que faz com que os estudantes tenham
dificuldade em aplicar o conhecimento para resolver problemas reais, que são
sempre complexos e misturados. Embora esse modelo tenha ajudado a aprofundar os
estudos em áreas específicas no passado, hoje ele prejudica a formação prática.
No desenvolvimento de software, criar um aplicativo exige, ao mesmo tempo,
conhecimentos de lógica de programação, facilidade de uso, regras de negócio e
infraestrutura \cite{paulista2018desenvolvimento}.

Para resolver essa divisão, a literatura recomenda integrar as matérias, com
destaque para os conceitos de multidisciplinaridade e interdisciplinaridade. A
multidisciplinaridade acontece quando várias disciplinas trabalham em um mesmo
projeto lado a lado, mas cada uma mantém seu próprio foco. De acordo com
\cite{bentes2024integracao}, usar uma visão multidisciplinar é o que faz a
diferença na hora de criar soluções integradas para problemas complexos. Na
criação de um produto digital, equipes multidisciplinares são necessárias para
cobrir todas as partes do sistema, mesmo que às vezes enfrentam dificuldades de
comunicação por não compartilharem a mesma "língua" técnica
\cite{silva2024processo}.

Por outro lado, a interdisciplinaridade vai além e cria uma união real e
contínua entre os conhecimentos. Como definem
\citeonline[p.~10]{scherer2011interdisciplinaridade}, a interdisciplinaridade é
uma "maneira complexa de entendimento e enfrentamento dos problemas" que
integra as disciplinas através do diálogo e da negociação constante. Na prática
da sala de aula, isso significa que a escolha feita na arquitetura de um
sistema (uma disciplina) afeta diretamente a qualidade do código e a estrutura
do banco de dados (outras disciplinas). É dessa mistura que o ensino deixa de
ser apenas a transmissão de regras isoladas e se torna um espaço para o
desenvolvimento completo do aluno, preparando-o para o mercado.

Para que essa integração de matérias funcione, é preciso mudar a forma de
ensinar e de aprender. O modelo tradicional, no qual o professor apenas fala e
o aluno escuta, não serve para resolver problemas práticos que exigem
conhecimentos variados \cite{ribeiro2025interdisciplinaridade}. É aí que entram
as Metodologias Ativas, com destaque para a Aprendizagem Baseada em Problemas
(PBL) \cite{ribeiro2025interdisciplinaridade}. Essas metodologias mudam a
rotina comum da sala de aula. Para
\citeonline{ribeiro2025interdisciplinaridade}, o uso de metodologias ativas
"favorecem aprendizagens profundas e competências amplas". No ensino de
tecnologia, isso imita os desafios reais das empresas de software. Quando os
alunos recebem um problema aberto, como criar um sistema para organizar as
demandas de um cliente real, eles precisam buscar os conhecimentos de várias
matérias simultâneas para conseguir criar uma solução que funcione.

Esse esforço gera o protagonismo do estudante, que é o fator principal para um
aprendizado de qualidade, pois faz o aluno deixar de ser um espectador e passar
a ser o centro do próprio aprendizado \cite{ribeiro2025interdisciplinaridade}.
O aluno protagonista entende o motivo de estar estudando aquele assunto, pois
vê a teoria sendo aplicada na hora. Como reforça
\citeonline{bentes2024integracao}, as universidades têm o dever de formar
profissionais que saibam se adaptar e que consigam usar o que aprenderam em
diferentes situações. Assim, ao invés de aprender apenas a usar uma ferramenta,
o estudante ganha a flexibilidade necessária para inovar e criar produtos de
software realmente úteis.

Embora a Aprendizagem Baseada em Problemas seja muito conhecida, existem outros
modelos e métodos de trabalho que se alinham perfeitamente com a criação de
sistemas. A Aprendizagem Baseada em Projetos, por exemplo, exige que os alunos
construam um produto real do início ao fim
\cite{ribeiro2025interdisciplinaridade}. No contexto tecnológico, isso
significa planejar, programar e entregar um software completo e funcional ao
longo de um semestre.

Para dar conta dessa entrega, as metodologias ágeis são trazidas do mercado de
trabalho direto para a sala de aula. Os alunos se organizam em equipes de
desenvolvimento, dividem as tarefas em ciclos de entrega e gerenciam o próprio
tempo para alcançar o objetivo final do projeto. Como define
\citeonline{sutherland2019scrum} em sua obra "Scrum: a arte de fazer o dobro do
trabalho na metade do tempo", esse formato ágil permite que as equipes resolvam
problemas complexos de forma mais rápida e com foco na qualidade contínua, o
que se adapta muito bem ao ritmo e aos prazos do ambiente acadêmico.

Dentro dessas equipes, o desenvolvimento do software ocorre por meio da
aprendizagem cooperativa. De acordo com \citeonline{johnson1994learning}, para
que a cooperação exista de fato, é necessário que o sucesso do projeto dependa
do esforço mútuo e da responsabilidade de todos. Na prática, cada estudante
assume uma tarefa técnica no sistema, mas todos precisam dialogar
constantemente para que as partes do código se conectem sem erros. Quando
surgem bloqueios técnicos durante essa construção, aplica-se naturalmente a
Instrução aos Pares (Peer Instruction). Criada pelo professor
\citeonline{mazur1996peer} na década de 1990, essa abordagem foca na
participação ativa, incentivando que os próprios alunos debatam os conteúdos e
encontrem as respostas juntos. Dessa forma, um aluno que tem mais domínio sobre
a infraestrutura pode explicar um conceito para um colega focado no visual do
sistema, fortalecendo a compreensão de toda a equipe.

Além da organização humana, a assertividade técnica da solução depende de
métodos práticos de engenharia, destacando-se a prototipagem rápida e os ciclos
curtos de revisão. Antes de gastar tempo escrevendo o código definitivo, os
estudantes criam simulações visuais do sistema.
\citeonline{coutinho2006prototipagem} explica que a prototipagem rápida
funciona como uma etapa fundamental no desenvolvimento ágil para avaliar
funcionalidades complexas e validar ideias de design de forma barata e
antecipada. Trabalhando em conjunto com a prototipagem, as equipes fazem
pequenas entregas parciais para receber avaliações (feedbacks) frequentes de
professores e colegas. Esse processo contínuo garante que o grupo consiga
ajustar a lógica do software e corrigir o rumo do projeto muito antes da
apresentação final.

\section{Eixos de Integração no Desenvolvimento de Software: O papel dos Projetos Integrados (PIES)}\label{sec:eixos-de-integracao}

Os Projetos Integrados em Engenharia de Software (PIES) funcionam como o
coração do currículo, pois é neles que o aluno deixa de ver as matérias de
forma isolada e passa a aplicá-las na construção de um sistema real.
\citeonline{scherer2011interdisciplinaridade} apontam que a união de diferentes
conhecimentos ajuda o estudante a enxergar o aprendizado de um jeito mais amplo
e serve para resolver problemas práticos do dia a dia. No contexto do curso,
essa estratégia serve para quebrar a ideia de que cada disciplina é uma
"caixinha" fechada e sem conexão com as outras. Essa integração acontece de
forma progressiva ao longo do curso, dividida em três grandes etapas que
aumentam de complexidade para simular os desafios reais do mercado de trabalho.

\subsection{ PIES 1: O Alicerce do Desenvolvimento (Dados e Projeto)} \label{sec:pies-1}

No primeiro estágio do projeto integrado, os alunos começam a definir o escopo
do sistema que vão criar, focado em construir a base técnica e a estrutura
inicial do software. Nesse período, as disciplinas de Projeto Detalhado de
Software e Fundamentos de Bancos de Dados trabalham juntas e dão suporte uma à
outra. Na matéria de Projeto Detalhado, o estudante aprende regras essenciais
para organizar o código, como a separação de interesses, o encapsulamento de
informações, a coesão e o acoplamento. Essas regras garantem que o código seja
limpo e fácil de mexer no futuro. Além disso, o aluno precisa lidar com
questões práticas como o tratamento de exceções (erros no sistema), a
persistência dos dados e a programação orientada por responsabilidade,
desenhando componentes e criando interfaces seguras entre os módulos do sistema
por meio de padrões de projeto (design patterns).

Ao mesmo tempo, todo esse desenho do software precisa ser guardado em algum
lugar, e é aí que entra a disciplina de Fundamentos de Bancos de Dados. O aluno
aprende a arquitetura de um Sistema Gerenciador de Banco de Dados (SGBD) e
estuda como fazer a modelagem usando o Modelo Entidade-Relacionamento e o
Modelo Relacional. Na prática do PIES 1, o estudante traduz os diagramas de
classes que fez no projeto detalhado em tabelas reais no banco de dados,
escrevendo comandos em SQL para inserir, buscar e alterar as informações. Para
garantir que o sistema não tenha dados duplicados ou quebrados, aplicam-se
conceitos avançados de restrições de integridade e normalização de tabelas.

Como mostram \citeonline{paulista2018desenvolvimento}, quando o aluno usa a lógica de programação
junto com a modelagem de dados para resolver um problema, ele consegue
resultados muito mais precisos e corretos. A vivência prática no PIES 1 faz o
estudante entender que um bom código não funciona sozinho se o banco de dados
não estiver bem estruturado. Assim, a união dessas matérias evita erros logo no
início do desenvolvimento, garantindo que o sistema comece com uma base firme
para continuar crescendo nos semestres seguintes.

\subsection{ PIES 2: Expansão para a Web e Visão de Negócio} \label{sec:pies-2}

No segundo estágio, o projeto sai do computador local e ganha uma utilidade
prática na internet, sendo preparado diretamente para o mercado de trabalho.
Esse bloco é formado pelas disciplinas de Desenvolvimento de Software para Web,
Verificação e Validação de Software \(V\&V\) e Empreendedorismo. O papel do
Desenvolvimento Web é fazer o sistema funcionar no navegador. Para isso, os
alunos aprendem sobre a internet e seus protocolos de comunicação, além de
dominarem o uso de HTML e CSS para criar a estrutura e o visual das páginas. A
disciplina também exige o uso de design responsivo e boas práticas de
acessibilidade para que o site funcione bem em qualquer tamanho de tela. Com o
uso do JavaScript, os estudantes aplicam lógica no navegador e fazem chamadas
assíncronas, permitindo que o sistema atualize as informações na tela sem
precisar recarregar a página inteira.

Para garantir que esse sistema web funcione sem falhas e com segurança, os
alunos aplicam ao projeto as técnicas da disciplina de Verificação e Validação
\(V\&V\). Eles aprendem a fazer o planejamento e a documentação das estratégias
de teste ao longo de todo o ciclo de vida do produto. Na prática, a equipe
realiza a análise estática do código e faz revisões de software para achar
problemas antes mesmo de rodar o programa. Em seguida, criam testes de unidade,
analisam a cobertura do código e usam técnicas de teste funcional (caixa preta)
e testes de integração, garantindo que a união do front-end com o banco de
dados dê certo. Os casos de teste são criados com base nas histórias de
usuários e casos de uso desenvolvidos no projeto. O grupo também usa
ferramentas de teste automatizadas ligadas a sistemas de integração contínua,
analisa relatórios de falha, aprende a isolar defeitos por meio da depuração
(debug) e faz o acompanhamento dos problemas usando ferramentas de rastreamento
(tracking), seguindo padrões conhecidos de qualidade.

Essa maturidade técnica ganha um sentido comercial por meio do
Empreendedorismo. Os alunos aprendem sobre o perfil do empreendedor, o
desenvolvimento de habilidades de negócios e a importância das pequenas
empresas na economia. A ementa aborda temas como inovação, propriedade
intelectual, transferência de tecnologia e o papel de incubadoras e parques
tecnológicos no suporte a novas ideias. Os estudantes utilizam ferramentas de
plano de negócios para entender o mercado, avaliar fontes de financiamento e
estruturar a viabilidade financeira do software que estão criando no PIES 2. De
acordo com \citeonline{bentes2024integracao}, as universidades funcionam como
motores de inovação quando conseguem conectar o que é estudado em sala de aula
com o que a sociedade e o mercado realmente precisam, diminuindo a distância
entre teoria e a prática. Desse modo, os testes de qualidade e o
desenvolvimento web passam a valorizar o plano de negócios do grupo,
transformando um exercício escolar em um produto com real potencial de mercado.

\subsection{ PIES 3: Engenharia Avançada e Qualidade de Software} \label{sec:pies-3}

Na fase final do ciclo de projetos integrados, o nível de exigência sobe para o
padrão das grandes empresas de tecnologia, onde o foco vai além de apenas fazer
o sistema funcionar: é preciso que ele seja seguro, aguente muitos acessos,
seja fácil de gerenciar e agrade ao usuário. Esse eixo maduro reúne as
disciplinas de Arquitetura de Software, Desenvolvimento para Dispositivos
Móveis, Gerência de Projetos, User Experience (UX) e Qualidade de Software. Na
matéria de Arquitetura de Software, o aluno aprende a definir a estrutura de
alto nível do sistema e estuda diferentes estilos arquiteturais, como
arquitetura em camadas, baseada em eventos, cliente-servidor ou filtros e
tubulações (pipes-and-filters). O estudante analisa a relação de custo e
benefício entre os atributos da arquitetura, cuida de questões de hardware,
garante a reutilização de código e faz a rastreabilidade dos requisitos para
que a estrutura do software respeite o que o cliente pediu.

Essa arquitetura robusta serve de base para a disciplina de Desenvolvimento de
Software para Dispositivos Móveis. Os alunos estudam os desafios da computação
móvel e aprendem a programar para plataformas como o Android SDK, criando
interfaces e layouts específicos para celulares e tablets. Eles aprendem a
trabalhar com ciclos de vida de telas (activities e intents), serviços em
segundo plano (services), mapas, localização e o uso de sensores disponíveis
nos aparelhos, como o acelerômetro e o GPS. Para que o aplicativo seja
realmente útil e fácil de usar, aplicam-se os conceitos de User Experience
(UX). Os estudantes aprendem métodos e técnicas para conhecer o comportamento
do usuário, aplicando o design de serviços no ambiente digital para criar
fluxos de navegação que sejam agradáveis e evitem a frustração de quem está
usando o aplicativo.

Para organizar todo esse trabalho técnico e de design, a Gerência de Projetos
traz os conceitos de ciclo de vida, gestão de interessados (stakeholders) e
estratégias de seleção. Os alunos gerenciam o escopo, controlam os custos
(orçamentos) e criam cronogramas gerenciando o tempo por meio da definição e
sequenciamento de atividades. Eles aplicam metodologias ágeis de mercado, como
o Scrum \cite{sutherland2019scrum}, dividindo as tarefas em ciclos de entrega e
controlando os riscos e as comunicações da equipe de forma organizada.

Por fim, a matéria de Qualidade de Software define as métricas e os padrões que
o produto deve atingir, usando modelos conhecidos como a ISO 9126-1, o CMMI e o
MPS.BR. Os alunos fazem o planejamento e a garantia da qualidade, realizam
auditorias e inspeções no código e analisam as causas de defeitos para evitar
que eles aconteçam novamente, desenvolvendo planos de qualidade oficiais.
\citeonline{silva2024processo} defendem que o desenvolvimento de produtos
digitais exige que diferentes visões trabalhem juntas para que o resultado
final seja realmente bom para o usuário. O PIES 3 consolida essa união de
gerência, arquitetura, desenvolvimento móvel, UX e qualidade, garantindo que o
software final funcione perfeitamente tanto no computador quanto no celular,
sem retrabalho e com o alto padrão exigido pelo mercado de tecnologia.
